% BEGIN RCB:R51_TITLE
\paragraph{Result R51: the constructed Deligne comparison maps and constituent coupling.}
% END RCB:R51_TITLE

% BEGIN RCB:R52_CONTINUATION
The full connection in the same cyclic coefficient frame is
now (XD.u1)--(XD.u7).  Define its columns by the literal monic
division
\[
 \Phi_t(S)S^b=(\chi(S)-t)Q_b^\Phi+R_b^\Phi,\qquad
 (C_\chi(t))_b=R_b^\Phi,\qquad (B_\chi)_b=(Q_b^\Phi)' .
\]
The differential remainder is \(R_b^\Phi-u(Q_b^\Phi)'\).
With the original branch and all oriented contours retained,
\[
 \Omega_u=-C_\chi(t)/u^2+B_\chi/u,\qquad
 \Pi_u=\Pi\Omega_u,\qquad
 \Omega_t=-A(t)/u,\qquad \Pi_t=\Pi\Omega_t,
 \qquad
 \partial_u\Omega_t-\partial_t\Omega_u+
                          [\Omega_u,\Omega_t]=0 .
\]
The degree-zero connection is
\(\partial_u-\Phi_t/u^2+1/u\); the degree-one connection is
\(\partial_u-\Phi_t/u^2\).  Their cochain identity, complete
division, and locally uniform contour differentiation are proved
in (XD.u1)--(XD.u6).  Here
\[
 C_\chi(t)=\Phi_t(A(t)),\quad
 (B_\chi)_{bb}=(b+1)/(q+1),\quad
 \partial_tB_\chi=0,\quad \operatorname{Tr}B_\chi=q/2.
\]
Thus the preceding determinant retains its full matrix
connection, with the same original constant
\[
 K_q=(q+1)^{-q/2}e^{\pi iq/2}\det W
          \prod_{b=0}^{q-1}\Gamma((b+1)/(q+1)),\qquad
 W_{jb}=e^{2\pi i j(b+1)/(q+1)}-1.
\]

On the exact primary slice \(\chi=h=(s-\rho)^m\), \(t=0\),
set \(n=m+1\), \(c=-(-\rho)^n/n\),
\(\beta_b=(b+1)/n\) and
\(d_b=n^{\beta_b-1}e^{\pi i\beta_b}\Gamma(\beta_b)\).
The original Pascal matrix has
\(P_{ab}=\binom ba\rho^{b-a}\) for \(a\le b\), and
the ordered period matrix is exactly
\(\Pi_s=e^{c/u}WD(u)P\), where
\(D(u)=\operatorname{diag}(d_bu^{\beta_b})\).
For an original \(p\)-dimensional ideal inclusion
\(I=J_{m-p}C_F\), whose \(J_{m-p}\) columns are
\((s-\rho)^{m-p},\ldots,(s-\rho)^{m-1}\) in the unchanged
\(s\)-basis, (SC.u5)--(SC.u7) compute
\[
 \det(I^*\Pi_s^*\Pi_s I)
 =|\det C_F|^2n^p(p+1)e^{2p\operatorname{Re}(c/u)}
                \prod_{b=m-p}^{m-1}|d_b|^2|u|^{2\beta_b}.
\]
The full Gram is
\(e^{2\operatorname{Re}(c/u)}
C_F^*D_{m-p}^*(W_{m-p}^*W_{m-p})D_{m-p}C_F\), with
\(W_{m-p}^*W_{m-p}=n(I_p+\mathbf1\mathbf1^*)\).
Its comparison with the original theta Gram divides the
determinant by \(\det(I^*G_NI)\), and uses
\((I^*G_NI)^{-1}\) on that complete matrix.
The proof retains the original source unit and every frame
determinant.  Its image plane is constant in \(u\), so its normal
\(u\)-map and \(u\)-curvature vanish on precisely this primary
\(t=0\) slice.  The general two-parameter normal map remains
\(N\Pi\Omega_a I\) for \(a=u,t\), as in (SC.u1)--(SC.u4).

For AP's original packet polynomial \(h\), full fixed Taylor
representative \(\nu\), and fixed poles \(p\in P_\nu\), the
localized complex now has the entire connection
(APU1)--(APU23).  Put \(f=\Phi_h-ts\) and
\[
 Q_p=\frac{f(s,t)-f(p,t)}{s-p}-\frac{h(s)-t}{\deg h+1}.
\]
Its ordinary quotient is the constant \(1/(\deg h+1)\), so
its derivative correction is exactly zero.  In AP6's original
polynomial and simple-pole columns,
\[
 \Omega_u^{\rm loc}
 =\begin{bmatrix}
      \Omega_u^{\rm pol}&-Q/u^2\\
      0&-\operatorname{diag}(f(p,t))/u^2
   \end{bmatrix}.
\]
The weighted-residue matrix is
\(\mathcal R=[0,W_{\rm res}]\), with
\(W_{\rm res}=\operatorname{diag}(e^{f(p,t)/u})\);
the positively oriented loop map is exactly
\((2\pi i)\mathcal R\).
Both maps are horizontal in \(u,t\).
The same complete \(\nu\) intertwines these connections on
the transported gauge frame, with target weight \(w/\nu\).
The original endpoint-source map is
\(\sigma_h[\nu]J_0\), with every inverse Taylor coefficient
retained, as proved in (APU18)--(APU22).
Finally the original entire transport
\(C_a=\Pi_{\rm loc}(u,t)^{-1}\Pi_{\rm loc}(u,t+ua)\)
has the exact derivative
\[
 (C_a)_u=-\Omega_u^{\rm loc}(u,t)C_a+
 C_a\bigl(\Omega_u^{\rm loc}(u,t+ua)
                   +a\Omega_t^{\rm loc}(u,t+ua)\bigr).
\]
Its \(t\)-derivative and the differentiated weighted-residue and
loop diagrams are proved in (APU23).  The extra
\(a\Omega_t^{\rm loc}\) is the derivative of the original
parameter shift.  The identification \(h=\chi_{h,1}\) is the
literal variable isomorphism (PC4); at higher tensor degree the
cyclic and product families retain their specified
\(\eta_k\) and source maps rather than being assigned a common
coefficient rank.
% END RCB:R52_CONTINUATION

% BEGIN RCB:R56_CONTINUATION
The same weighted observer now has its full singular-parameter
equation (PC.u1)--(PC.u3):
\[
 \mathcal C_{k,u,t}=\mathcal C_k\Pi_\chi(u,t)^{-1},\qquad
 \partial_u\mathcal C_{k,u,t}
   =\mathcal C_k(C_\chi(t)/u^2-B_\chi/u)\Pi_\chi^{-1},
 \qquad
 \ker\mathcal C_{k,u,t}=\Pi_\chi\ker\mathcal C_k.
\]
Here the polynomial is literally \(\chi=\chi_{h,k}\);
the product-period restriction on \(E_h^{\otimes k}\) is
instead the actual rectangular form
\(\eta_k^*\Pi_{\rm prod}^*\Pi_{\rm prod}\eta_k\).
The source identity and parameter cyclic relation following
(XD.u7) are the exact maps between those two domains.
For the original kernel inclusion \(I_k\), its complete
normal maps are \(Z_a=N\Pi_\chi\Omega_a I_k\), \(a=u,t\),
and all four curvature coefficients are
\[
 \partial_{\bar a}\partial_b\log\det H
       =\operatorname{Tr}(H^{-1}Z_a^*Z_b).
\]
At \(t=0\), the order-two coefficient preserves this ideal,
so \(Z_u=u^{-1}N\Pi_\chi B_\chi I_k\), \(Z_t=0\).
These are calculated matrix entries on the same kernel; neither
the full ideal nor its nilpotent jets have been suppressed.

The actual tensor coefficient map is also now retained in
(PC30a)--(PC30g).  On an ordered local tuple let
\(x_i=s_i-\rho_i\), \(\lambda=\sum_i\rho_i\),
\(D=\sum_i(m_i-1)\), and
\(U_{\boldsymbol\rho}=\prod_i j_{\rho_i}(\upsilon_h)\).
Over the specified coefficient localization
\(\mathscr R=A_{\ker\alpha}\), for an actual map to characteristic
\(p>\max_i m_i\), the exact map and its saturation basis are
\[
 V_j=\sum_{\substack{0\le b_i<m_i\\|\mathbf b|=j}}
            \frac{x^{\mathbf b}}{\prod_i b_i!},\quad
 W_j=U_{\boldsymbol\rho}V_j,\quad
 \eta_{\boldsymbol\rho}((S-\lambda)^j)=j!W_j,
 \qquad
 \Lambda_U=\bigoplus_{j=0}^D\mathscr R W_j.
\]
Its full containing tensor module retains the explicitly split
free complement, and
\[
 \Lambda_U/\eta_{\boldsymbol\rho}\mathcal C_{\mathscr R}
    \simeq\bigoplus_{j=0}^D
                  \mathscr R/(p^{v_p(j!)}),\qquad
 \ker\overline\eta_{\boldsymbol\rho}
      =((S-\overline\lambda)^p),\qquad
 \operatorname{im}\overline\eta_{\boldsymbol\rho}
      =\bigoplus_{j=0}^{\min(D,p-1)}B\overline W_j .
\]
Here \(B\) is the specified target \(k_0\), or the explicit
quotient \(\mathscr R/p\mathscr R\). The kernel is zero when \(D<p\).
The original source index is \(D+1\), while the index on the
actual specialized image is \(\min(p,D+1)\); the full free
resolution and Tor connecting map in (PC30e) prove their exact
relation.  When the specified map includes all one-factor CRT
data and \(p>m_\rho\) for every retained primary order, (PC30f)--(PC30g) also retain every new residual collision:
\[
 \psi_{h,k}(X)=
 \prod_{\mu}
 (X-\mu)^{
 \max_{\alpha(\sum_i\rho_i)=\mu}
            \min(p,1+\sum_i(m_i-1))},\quad
 \ker\overline\eta_k=(\psi_{h,k})/(\overline\chi_{h,k}),\quad
 \operatorname{im}\overline\eta_k\simeq k_0[X]/(\psi_{h,k}).
\]
The product is over attained residual sums; if the carrier is empty,
its product is \(1\) and both modules are zero.
The last arrow is the proved module map
\([P]\mapsto\overline U_kP(\overline S)\); its unweighted
counterpart is the corresponding algebra map.

For the complete primary packet at \(t_i=0\), the ordered
product boundary is explicitly calculated in
(SPBT.1)--(SPBT.3), (BPC.1)--(BPC.14),
(BTP.1)--(BTP.11), and (BTF.1)--(BTF.11).
Writing \(n=m+1\), its ordinary complex invariant dimension is
\[
 d_{0,k}=\frac{m^k+m(-1)^k}{m+1}.
\]
With \(E_0\) selecting \(\sum_i(b_i+1)\equiv0\pmod n\),
the exact coefficient, source and period projectors satisfy
\[
 Q_0=P_k^{-1}E_0P_k,\quad
 \widehat Q_0=U_kQ_0U_k^{-1},\qquad
 \mathcal P_0\Pi_{\rm prod}U_k^{-1}
       =\Pi_{\rm prod}U_k^{-1}\widehat Q_0.
\]
The terminal original monomial survives precisely when
\(n\mid k\), and retains its arithmetic eigenvalue
\(a^{k\rho}\).  The exact compression and its outgoing
nilpotent map are (BTP.7)--(BTP.8).
On the actual phase specialization with \(p>m+1\), the
common additive factor remains \(\mathcal L_\psi(kc/u)\):
the invariant dimension is \(d_{0,k}\) when \(kc=0\) and
zero otherwise, and the Swan conductor is respectively zero
or \(m^k\).  Every surviving boundary orbit retains its full
product of negative Gauss constants and weight \(k\).
The common cyclotomic projector maps in (BTF.6) give their
exact two base changes; they do not assign a characteristic-\(p\)
coefficient scalar to a characteristic-zero Gauss line.
The amplified coefficient still contains its original \(K!\),
\(K=k(m-1)\), with image zero exactly when \(K\ge p\)
on a coefficient stage retaining the inverse unit.
The case \(m=1\) has \(K=0\) and no factorial loss.
Moreover, in the retained off-critical quartet CRT corner
for \(k\ge p\), (BTF.9) proves
\[
 1=p(\rho+\overline\rho-1)\mathscr D,\qquad
 p^{-1}=(\rho+\overline\rho-1)\mathscr D
\]
from the actual separating idempotent's divided difference.
That larger stage has no unital characteristic-\(p\) map;
the original local coefficient map, its precise saturation,
kernel, Tor class and character remain as displayed above.
% END RCB:R56_CONTINUATION

% BEGIN RCB:R57_CONTINUATION
The same numerical-range calculation is now transported in
both original parameters by (LC.u1), with the unchanged
\(E_N(t)\) and \(G_N\):
\[
 \begin{aligned}
 \partial_u\widetilde G_N
   &=-\Pi^{-*}G_N\Omega_u\Pi^{-1},&
 \partial_{\bar u}\widetilde G_N
   &=-\Pi^{-*}\Omega_u^*G_N\Pi^{-1},\\
 \partial_u\widetilde L_t
   &=\Pi[\Omega_u,L_t]\Pi^{-1},&
 \partial_u\widetilde A_t
   &=\Pi[\Omega_u,A_t]\Pi^{-1}.
 \end{aligned}
\]
Here \(\Omega_u=-C_\chi(t)/u^2+B_\chi/u\) is the full
original cyclic connection.  The exact original source
commutator retains
\(r_N[A,\Omega_u]j_E+dK_N^{\rm prim}\Omega_u j_E\),
as proved in (DS74) and carried through (LC.u1).
Thus the preceding original-metric parabola holds for every
retained \(u\ne0\), with its same control value and the
displayed moving metric, rather than an uncomputed target
norm.
The comparison numbers in (LC14)--(LC16) are now the actual
\(m(u,t),M(u,t)\) of \(G_N^{-1}\Pi(u,t)^*\Pi(u,t)\).
For each constituent and each neighboring inclusion,
(SC.u1)--(SC.u4) give the complete Hermitian Hessian
\(\operatorname{Tr}(H^{-1}Z_a^*Z_b)\);
at nonzero \(t\), both terms of \(\Omega_u\) and their
cross pairing remain.  On the precise single-primary ideal
slice the full matrix and ideal determinant are evaluated
by (SC.u5)--(SC.u7), including the original comparison Gram.
% END RCB:R57_CONTINUATION

% BEGIN RCB:R59_CONTINUATION
The entire same diamond now has its exact \(u\)-equation
(DC31), for all four formal roles \(s\in\{a,b,d,l\}\):
\[
 T_{s,u,t}=T_s\Pi(u,t)^{-1},\qquad
 \partial_uT_{s,u,t}
   =T_s(C_\chi(t)/u^2-B_\chi/u)\Pi(u,t)^{-1}.
\]
Its kernels are \(\Pi K_s\), and its specified right
inverses, intersection and sum diagrams are the same ones
in (DC23)--(DC30).  Multiplication by the original common
quotient maps proves their differential compatibility,
retaining all four roles and every inverse-unit factor.
For the actual circle weights put
\(D_\nu=\operatorname{diag}(\nu_n)_{n\in\mathcal Z_L}\).
The exact observed form and its derivatives are (DC32):
\[
 \begin{aligned}
 H_b^{\rm obs}&=\Pi^{-*}T_b^*D_\nu T_b\Pi^{-1},
 &\ker H_b^{\rm obs}&=\Pi K_b,\\
 \partial_uH_b^{\rm obs}
   &=-\Pi^{-*}T_b^*D_\nu T_b\Omega_u\Pi^{-1},&
 \partial_{\bar u}H_b^{\rm obs}
   &=-\Pi^{-*}\Omega_u^*T_b^*D_\nu T_b\Pi^{-1}.
 \end{aligned}
\]
Every retained atom has its original positive weight; an
empty sample set gives the empty observation and its full
kernel.  The shared and additional circle sums, polynomial
source lift and receiving support labels therefore retain
their original maps while the parameter varies.
% END RCB:R59_CONTINUATION

% BEGIN RCB:R62_TRANSLATION_REPLACEMENT
The exact Pascal translation preserves the \(t\)-derivatives
of order at least two at fixed \(u\), while its scalar factor
retains the first-derivative trace shift
\(2p\operatorname{Re}a\).  Its pure \(u\)-derivatives are
the explicitly retained terms in (RCX.u5)--(RCX.u6) below.
% END RCB:R62_TRANSLATION_REPLACEMENT

% BEGIN RCB:R62_CONTINUATION
The same constant quotient lift and the complete two-parameter
period matrix give the exact compensating Hessian
\[
 \partial_{\bar a}\partial_b\log\det K
    =-\operatorname{Tr}(H^{-1}Z_a^*Z_b),\qquad
 a,b\in\{u,t\},\quad
 H=I^*\Pi^*\Pi I,\quad
 Z_a=N\Pi\Omega_a I .
\]
Indeed the displayed Schur determinant identity holds at
each \((u,t)\), its constant frame factor has zero
derivative, and \(\log|\det\Pi|^2\) is pluriharmonic
where \(\Pi\) is invertible.  Applying (SC.u2) proves
the equality for both diagonal and both mixed entries,
with the original quotient metric and frame preserved.
The earlier quadratic vanishing statement is its \(t\)
entry at \(t=0\); the \(u\) entry retains the
\(B_\chi/u\) normal term of (SC.u3).

The complete ordered derivative recursion now is
(RCX.u1)--(RCX.u2):
\[
 Q_{0,0}=I,\qquad
 Q_{r+1,s}=\partial_uQ_{r,s}+\Omega_uQ_{r,s},\qquad
 Q_{r,s+1}=\partial_tQ_{r,s}+\Omega_tQ_{r,s},\qquad
 \partial_u^r\partial_t^s\Pi=\Pi Q_{r,s}.
\]
Flatness proves compatibility of all derivative orders.
In particular,
\[
 \begin{aligned}
 Q_{2,0}
   &=C_\chi^2/u^4+
       (2C_\chi-C_\chi B_\chi-B_\chi C_\chi)/u^3
                       +(B_\chi^2-B_\chi)/u^2,\\
 Q_{1,1}
   &=C_\chi A_t/u^3+(A_t-B_\chi A_t)/u^2.
 \end{aligned}
\]
All earlier pure \(t\) jets, including \(3R^2/u^2\)
in the fourth derivative, are the same \(r=0\) slice.
Along a specified curve \((u(s),t(s))\), the exact
metric variation inserted into the original residue
Rayleigh-quotient formula is
\[
 \dot H_{\rm full}
   =\Omega_{\rm path}^*H_{\rm full}
            +H_{\rm full}\Omega_{\rm path},\qquad
 \Omega_{\rm path}=\dot u\,\Omega_u+\dot t\,\Omega_t,
 \quad H_{\rm full}=\Pi^*\Pi .
\]
The full forced-equation descendant also remains, by (RCX.u4):
\[
 2u\Pi_{tt}+u^2\Pi_{utt}+k\Pi_t+ku\Pi_{ut}
       +\Pi R+u\Pi_uR=\Pi\Omega_u L_t.
\]
It composes with the unchanged theta-source map \(j_E\)
and retains its original boundary primitive.

Finally, with \(p=\dim F\) and
\(h_a(t)=-\Phi(-a)-ta\), the original translated phase gives
\[
 \psi_a-\psi=2p\operatorname{Re}(h_a(t)/u),\qquad
 \partial_u(\psi_a-\psi)=-p h_a(t)/u^2,\qquad
 \partial_u^2(\psi_a-\psi)=2p h_a(t)/u^3,\qquad
 \partial_u\partial_t(\psi_a-\psi)=pa/u^2 .
\]
These are the exact derivatives proved in
(RCX.u5)--(RCX.u6), through the full translated
connection (DT.14)--(DT.16).  Their Hermitian mixed
derivatives vanish because the retained scalar
difference is pluriharmonic, while these pure
derivatives remain in the ordered recursion.
% END RCB:R62_CONTINUATION
